\chapter{Conclusión}

El presente Seminario de Título demuestra que el diseño de un agente basado en modelos de lenguaje de gran escala para el análisis semántico de incidentes de seguridad constituye una propuesta técnicamente fundamentada y arquitectónicamente viable. A través de una revisión sistemática de 52~publicaciones científicas del período 2023--2026, seleccionadas mediante la metodología PRISMA~2020 a partir de un corpus inicial de 357~registros en Scopus y Web of Science, se caracterizó el estado del arte en cuatro dimensiones: \textit{log parsing}, detección de anomalías, automatización SOC, estrategias de prompting, y se identificaron tres brechas estructurales que los enfoques existentes no resuelven de manera conjunta.

La primera brecha es la ausencia de razonamiento causal iterativo: la totalidad de los trabajos de análisis de logs revisados opera como \textit{pipeline} lineal pasivo que produce una clasificación o resumen sin formular hipótesis ni refinarlas activamente con nueva evidencia. Esta limitación es señalada explícitamente en el corpus como la capacidad más ausente en los \textit{frameworks} actuales. La segunda brecha es la falta de integración efectiva con inteligencia de amenazas estructurada: solo 3 de los 52~trabajos del corpus integran MITRE ATT\&CK en su pipeline, y ninguno puebla una base vectorial desde el repositorio STIX~2.1 completo. La tercera brecha es la carencia de narrativas de incidente completas y accionables: incluso los sistemas más avanzados del corpus, como el sistema multiagente de Hmimou et al.~\cite{hmimou2025multiagent}, que alcanza el 93.6\,\% de \textit{accuracy} y una reducción del 41.3\,\% de falsos positivos frente a un SIEM baseline, o el sistema de Ali et al.~\cite{ali2025apt}, con 99.6\,\% de \textit{accuracy} y FPR$=0.0\,\%$, entregan clasificaciones o predicciones individuales sin integrar secuencia temporal, tácticas MITRE, niveles de confianza por hallazgo ni trazabilidad del razonamiento en un único reporte accionable.

La arquitectura propuesta responde directamente a dichas brechas mediante un pipeline \textit{Plan-and-Execute} implementado con LangGraph, que separa explícitamente la fase de planificación de la de ejecución para superar la miopía del patrón ReAct tradicional. Esta decisión se fundamenta en los hallazgos del corpus sobre estrategias de prompting: LogPrompt~\cite{liu2024logprompt} demuestra que el razonamiento estructurado mejora entre un 62.4\,\% y un 380.7\,\% el desempeño de los LLMs sobre prompts simples, mientras que estudios comparativos muestran diferencias de hasta 20~puntos de F1 entre estrategias zero-shot y few-shot. El sistema integra 7~nodos (4~agentes LLM especializados y 3~herramientas determinísticas), 3~aristas condicionales que gobiernan el flujo según el estado del análisis, una base vectorial ChromaDB con umbrales de confianza calibrados, complementada con consultas externas a NVD y AbuseIPDB para enriquecer el diagnóstico con inteligencia de amenazas en tiempo real. La versión preliminar del prototipo implementada en este semestre confirma la viabilidad del flujo completo: desde la agregación determinística de logs heterogéneos hasta la generación de un reporte estructurado con hipótesis del incidente, nivel de confianza, comandos de mitigación validados y trazabilidad del razonamiento.

Como trabajo futuro, el Proyecto de Título concentrará el esfuerzo en tres tareas. Primero, la evaluación experimental rigurosa del prototipo sobre datasets con \textit{ground truth} (HDFS, BGL y UWF-ZeekDataFall22), donde se medirán métricas de precisión en la identificación de incidentes, coherencia de las narrativas generadas y desempeño comparativo frente a enfoques basados en reglas. Segundo, la integración del repositorio oficial MITRE ATT\&CK CTI en formato STIX~2.1 para poblar la base vectorial ChromaDB con conocimiento adversarial estructurado, fortaleciendo la capacidad del agente para mapear eventos a tácticas y técnicas conocidas de forma sistemática. Tercero, el refinamiento de los componentes del agente a partir de los resultados experimentales, con el propósito de mejorar la trazabilidad del razonamiento y la calidad de las narrativas generadas. Estas líneas de trabajo permitirán consolidar al agente como una herramienta de apoyo operacionalmente efectiva ante la creciente complejidad de las amenazas dirigidas a organizaciones sin capacidad de respuesta especializada.

\section{Cumplimiento de los Objetivos}

\section{Principales Contribuciones}

\section{Limitaciones}

\section{Trabajo Futuro}